\begin{thesisabstract}{finnish}

Suuria kielimalleja hyödyntävät monivaiheiset työnkulut
maksavat jokaisesta lukemastaan tokenista, joten
kontekstin tiivistäminen on luonteva keino karsia
kustannuksia. Tiivistämistä mitataan kuitenkin yhden
fragmentin kysymys--vastaus-tarkkuudella, vaikka
tuotannon työkuormat yhdistävät tietoa useasta erikseen
tiivistetystä fragmentista. Sitä, säilyttääkö yhden
fragmentin laadun säilyttävä tiivistäminen myös tämän
vaikeamman ominaisuuden, ei ollut aiemmin tutkittu.

Ei säilytä. Uudessa $150$ instanssin vertailuaineistossa
neljä ilman koulutusta toimivaa tiivistäjää kymmenen
tiivistyssuhteen yli ei tuota lainkaan positiivista
työkuormatason korrelaatiota lähdesäilyvyys-F1:n ja
koordinaation onnistumisen välille (Spearmanin
järjestyskorrelaatio $[-0{,}82, +0{,}32]$). Romahdus on
äkillinen: jyrkkä koordinaation jyrkänne $\tau^{*}$
esiintyy $11$ solussa $12$:sta, kun tiivistäminen karsii
ne kriittiset tokenit, jotka suunnittelijan on
yhdistettävä. Kumuloituvan virheen malli selittää
kynnyksen ja palauttaa sen sijainnin ensimmäisen
kertaluvun tarkkuudella. Jyrkänne on tiivistäjän ja
tehtävän, ei suunnittelijan, ominaisuus: kalibroidulla
alueella sen sijainti ei siirry paikallisesta
Llama-3.1-8B:stä huipputason Qwen-2.5-72B:hen.

Työn tuloksena on käytäntöjä valvova muistiväylä, joka
tiivistää jokaisen fragmentin matkalla ja muuntaa
tokensäästöt suoraan pienemmäksi päättelykustannukseksi.
Sen tiivistystä ohjaa mittauksiin perustuva
jyrkännetietoinen kääre (CAAC), joka perääntyy
aggressiivisimpaan suhteeseen, joka pitää kunkin
fragmentin jyrkänteen turvalliselle puolelle. Avoimen
lähdekoodin viitetoteutus on toistettavissa yhdellä
komennolla esilasketun tiivistysvälimuistin avulla.

\end{thesisabstract}
